\section{Main takeaways}
\begin{itemize}
  \item A two-dimensional powder retains radial and transverse detector information that separates specular-family, off-specular, wavelength-sensitive, and mosaic-sensitive observables more effectively than a conventional one-dimensional reduction.

  \item The forward model begins with continuous rod intensities $S_{hk}(L,\lambda_s)$, distributes them through the two-dimensional-powder and crystallite-orientation densities, restricts the resulting reciprocal-space density to each state-specific Ewald surface, and maps it to the calibrated detector with the same finite integration used for the data.

  \item In the Bi$_2$Se$_3$ $r=0$ region, the long $003$ extension is radial in $2\theta$ and is reproduced by the wavelength-resolved calculation, while the transverse $\phi$ width remains the cleaner crystallite-tilt constraint.

  \item The bright near-origin $r=0$ response motivates a separate Parratt optical branch, with N\'evot--Croce roughness, that is smoothly connected to the finite-stack kinematic branch. In the present work this construction distinguishes the near-critical response from the Bragg and mosaic contributions. It is not used as an independent reflectivity refinement.

  \item PbI$_2$ stacking faults act by changing trilayer pair correlations within the continuous rod intensity. Off-specular rods retain basal-plane phase contrast and develop diffuse between-peak intensity, whereas $r=0$ is contrast-matched to the included stacking events and can contain faults without displaying comparable diffuse contrast.

  \item The PbI$_2$ comparison should test a nested hierarchy of nearly ordered 2H, faulted 2H, 2H+6H, and 2H+6H+4H while keeping detector geometry, mosaicity, ordered baseline, projection, and background treatment fixed.

  \item The staged refinement order is essential: instrumental geometry is determined first, sample alignment second, wavelength and mosaic effects third, ordered rod intensity next, and stacking-disorder parameters last. This ordering limits compensation among parameters and keeps the structural interpretation tied to matched detector-space observables.

  \item Detector-space indexing establishes where allowed reflections should appear. The present contribution is the subsequent continuous intensity and profile calculation that combines those locations with source phase space, rod intensity, orientation density, optical weights, and finite detector integration.\cite{SmilgiesLi2026IndexGIXS}
\end{itemize}
